\section{Experiments}
\label{sec:experiments}

The current experiments compare checkpoint-free heuristic agents under one shared benchmark protocol. All methods are evaluated on the same frozen test split and seed-controlled episode stream. The table below summarizes retained heuristic baseline results over \(5000\) test episodes.

\begin{table}[t]
  \centering
  \caption{Heuristic benchmark comparison on the test split under the shared evaluation protocol.}
  \label{tab:core_benchmark_results}
  \small
  \begin{tabular}{lrrrr}
    \toprule
    Method & Avg Profit (USD) \(\uparrow\) & Deal Rate \(\uparrow\) & Walkaway Rate \(\downarrow\) & Avg Rounds \\
    \midrule
    Random & 6058.8570 & 0.5796 & 0.4204 & 1.3208 \\
    Concession & \textbf{14725.4262} & \textbf{0.7244} & \textbf{0.2756} & \textbf{1.7196} \\
    \bottomrule
  \end{tabular}
\end{table}

These results show that the environment separates weak random pricing from a simple structured negotiation heuristic. The concession policy achieves higher seller profit and a higher deal rate than random behavior, which indicates that the benchmark rewards coherent price adaptation rather than merely ending negotiations quickly.

The current result set is intentionally compact. It establishes a runnable benchmark artifact and heuristic reference points. The next experimental step is to add zero-shot LLM baselines under the same fixed episode stream and to report invalid-output rate, API cost, and token usage alongside negotiation metrics.
